
où la
signification des flèches est évidente, et dont il suffit de vérifier la commutativité pour les éléments \(1+Xt,1+Yt,1+Zt\) de \(\Ghat(\mathbb Z[X,Y,Z])\) (resp. \(1+Xt,1+Xt\) de \(\Ghat(\mathbb Z[X,Y])\)) d'après 1.10, vérification qui est immédiate.

D'après 1.8, le produit par cette loi de deux séries se calcule par substitution de leurs coefficients dans les polynômes universels de la loi. Le \(n\)-ième polynôme universel sera noté \(P_n\); c'est un polynôme de \(\mathbb Z[X_i,Y_j]_{1\leq i,j\leq n}\), isobare de poids \(n\) par rapport à chacun des deux groupes de variables (voir 1.15), et symétrique par rapport à ces deux groupes. Donnons à titre d'exemple la valeur de \(p_2\) :
\[
  P_2(X_1,X_2;Y_1,Y_2)=X_2Y_1^2+Y_2X_1^2-2X_2Y_2.
\]

On a également la relation suivante : soient \(K\in\operatorname{Ob}(\mathcal C_0)\), \(\psi\in\Ghat(K)\), \(a\in K\), alors
\begin{equation}\tag{2.3.2}
  \psi(t)\circ(1+at)=\psi(at).
\end{equation}
Comme d'habitude, elle se ramène par l'argument fonctoriel standard au cas \(K=\mathbb Z[X,Y]\), \(\psi=1+Xt\), \(a=Y\), où elle n'est autre que 2.3.1.

Appliquant 2.3.2 à \(\Ghat(K)=1+K[[t]]^+\), et à \(a=1\), on voit que l'élément \(1+t\) est unité pour cette multiplication. Pour éviter des confusions, nous le noterons parfois \(\underline{1}\), ou encore \(\underline{1}_\circ\).

Le foncteur \(\Ghat\) peut maintenant être considéré comme foncteur de \(\mathcal C_0\) dans la catégorie des anneaux commutatifs à éléments unité, et par conséquent être composé avec lui même~; pour éviter des confusions,
nous 
